%% Proofs for ED26-SCI-SaltelliAnnoni2010-001.
\begin{annexitem}{annex:p15.global.oat-obstruction}
 {thm:p15-global-oat-obstruction}{The interaction invisible on coordinate lines}
\begin{proof}
At an OAT point where coordinate \(1\) varies, \(x_2=b_2\).
Where coordinate \(2\) varies, \(x_1=b_1\).
On every other coordinate line both equalities hold.
Thus \(f_C=0=f_0\) on the whole design set, independently of \(C\).
The rectangle contains a point with both \(x_1\ne b_1\) and
\(x_2\ne b_2\); their product is nonzero.
At that point its magnitude grows without bound as \(|C|\) increases.
Both functions are continuous polynomials. Hence even this very simple
model class is not identified by all the OAT line values.
\end{proof}
\end{annexitem}

\begin{annexitem}{annex:p15.global.ball-bound}
 {thm:p15-global-ball-bound}{The inscribed ball and its exact scope}
\begin{proof}
Every centered OAT point differs from \(b\) in at most one coordinate,
by at most \(1/2\). Its Euclidean distance from \(b\) is therefore at
most \(1/2\).

Let \(v_d\) be the volume of the unit-radius ball in \(\mathbb R^d\).
Scaling coordinates gives volume \(v_dR^d\) at radius \(R\).
Using the last two coordinates as a disk and applying Tonelli,
\[
 v_{d+2}R^{d+2}
   =\int_{\|x\|\leq R}\pi(R^2-\|x\|^2)\,dx
   =\pi\int_0^{R^2}v_d(R^2-t)^{d/2}\,dt
   =\frac{2\pi}{d+2}v_dR^{d+2}.
\]
The middle equality writes the nonnegative height as an integral of
its superlevel sets. With \(v_1=2\) and \(v_2=\pi\), this recurrence
determines both parities.

The gamma recurrence, proved in
Theorem~\ref{thm:p15-gamma-normalization}, gives
\(\Gamma(s+1)=s\Gamma(s)\), \(\Gamma(1)=1\).
For the needed half-integer initial value, substitution \(t=x^2\)
gives
\(\Gamma(1/2)=\int_{\mathbb R}e^{-x^2}\,dx=:J\).
This integral is finite, since \(e^{-x^2}\leq e^{-|x|}\) for
\(|x|\geq1\). Tonelli and a disk-area superlevel calculation give
\[
 J^2=\int_{\mathbb R^2}e^{-\|z\|^2}\,dz
     =\int_0^\infty e^{-t}\pi t\,dt=\pi.
\]
Thus \(J=\sqrt\pi\), and induction in the ball recurrence gives
\(v_d=\pi^{d/2}/\Gamma(d/2+1)\).
Set \(R=1/2\); the cube volume is \(1\).
This gives every displayed ratio and the stated value at \(d=12\).
The recurrence
\(r_{d+2}=\pi r_d/(2(d+2))\)
has a ratio eventually bounded by \(1/2\), so both parity subsequences
tend to zero.

For \(d\geq2\), each coordinate line fixes at least one other coordinate
to a singleton of Lebesgue measure zero. Product integration makes its
\(d\)-dimensional measure zero. A finite union retains measure zero.
This last fact concerns the line set; it does not equate its measure
with the volume of the containing ball.
\end{proof}
\end{annexitem}

\begin{annexitem}{annex:p15.global.factorial}
 {thm:p15-global-factorial}{The finite corner basis}
\begin{proof}
Equip the corner set with uniform weight \(2^{-d}\).
For subsets \(u,v\),
\(\chi_u\chi_v=\chi_{u\triangle v}\).
If \(u\ne v\), at least one coordinate appears to its first power,
and summing over its two signs cancels. Thus
\[
 2^{-d}\sum_z\chi_u(z)\chi_v(z)
   =\begin{cases}1,&u=v,\\0,&u\ne v.\end{cases}
\]
There are \(2^d\) such orthonormal functions in the \(2^d\)-dimensional
space of all corner responses. They are a basis, proving the coefficient
formula, expansion and uniqueness.

For the main effect of \(i\), averaging at \(z_i=1\) and \(z_i=-1\)
cancels every term involving another coordinate and cancels the constant
between the two means. The difference is \(2\beta_{\{i\}}\).
At fixed \(z_j\), averaging over the remaining coordinates gives the
main \(i\)-effect
\(2\beta_{\{i\}}+2z_j\beta_{\{i,j\}}\).
Half its difference between \(z_j=1\) and \(z_j=-1\) is
\(2\beta_{\{i,j\}}\).
The general expansion also supplies all higher-order interactions.

The corner response space has eight coordinates when \(d=3\).
A fit with all eight basis functions interpolates those values exactly,
so there is no residual observation dimension.
Finally \(\prod_i x_i(1-x_i)\) is zero at all corners and equals
\(4^{-d}>0\) at the center. Adding an arbitrary multiple of it changes
the interior response without changing any fitted corner.
\end{proof}
\end{annexitem}

\begin{annexitem}{annex:p15.global.regression}
 {thm:p15-global-regression}{The linear projection and explained fraction}
\begin{proof}
The standardized variables have means zero and \(E[W^2]=1\).
For any coefficient vector \(b\),
\[
 E[(W-b^\mathsf TZ)^2]=1-2b^\mathsf Tc+b^\mathsf T\Sigma b.
\]
When \(\Sigma\) is positive definite, completion of the square gives
\[
 1-c^\mathsf T\Sigma^{-1}c
  +(b-\Sigma^{-1}c)^\mathsf T
           \Sigma(b-\Sigma^{-1}c).
\]
The minimum is unique at \(b=\Sigma^{-1}c\).
At that value \(c=\Sigma b\), so
\(R^2=b^\mathsf T\Sigma b=c^\mathsf T\Sigma^{-1}c\).
For uncorrelated standardized inputs, \(\Sigma=I\), yielding
the stated sum of squares.

For an exactly additive linear model, centering gives
\(Y-EY=\sum_i a_i\sigma_iZ_i\).
Independence gives
\(V=\sum_i a_i^2\sigma_i^2\).
The displayed coefficient vector therefore fits \(W\) exactly, proving
\(R^2=1\) and the coefficient formula.

For a finite sample, replace expectation by uniform averaging over its
observations and center and scale under that same averaging convention.
The entire quadratic calculation is unchanged, with empirical covariance
\(\widehat\Sigma\). A random sample from independent population variables
need not have zero empirical cross-products; the exact identity is then
\(R^2=b^\mathsf T\widehat\Sigma b\), not an unsupported sum of squares.
If a covariance matrix is singular, a null vector \(v\) has
\(E[(v^\mathsf TZ)^2]=0\), so adding it to a coefficient vector changes
no fitted response almost surely. Individual coefficients need not then
be unique, even though the minimizing projection is.
\end{proof}
\end{annexitem}

\begin{annexitem}{annex:p15.global.anova}
 {thm:p15-global-anova}{The product-law orthogonal decomposition}
\begin{proof}
Conditional expectation onto a coordinate sub-sigma-algebra can be
represented by integrating over the other coordinates under the product
law. Since \(Y\in L^2\), it is integrable, and iterated product integration
is justified by Fubini. In particular, for subsets \(u,v\),
\[
 P_uP_vY=P_{u\cap v}Y.
 \tag{1}
\]
To see this directly, first integrate the coordinates not in \(v\), then
those not in \(u\). A coordinate surviving both operations is exactly one
in \(u\cap v\); integrations of coordinates already removed contribute
their probability mass \(1\).
Conditional expectation is an \(L^2\) contraction, so all finite sums
defining the \(Y_u\) are square integrable.

The alternating-sum expression follows from the recursive definition by
finite subset inversion. Indeed, summing that expression over all \(u\)
gives coefficient
\(\sum_{u:v\subseteq u\subseteq I}(-1)^{|u|-|v|}
  =(1-1)^{d-|v|}\)
on \(P_vY\); only \(v=I\) survives, and \(P_IY=Y\).
The same calculation restricted to a subset gives
\(P_uY=\sum_{v\subseteq u}Y_v\), proving the recursive formula.

For \(i\in u\), apply \(P_{I\setminus\{i\}}\) to the alternating sum
for \(Y_u\). Pair the summand indexed by \(v\subseteq u\setminus\{i\}\)
with the one indexed by \(v\cup\{i\}\).
Identity (1) gives \(P_vY\) in both, with opposite signs.
Thus the conditional expectation is zero.

If \(u\ne v\), choose a coordinate in one of their set differences,
say \(i\in u\setminus v\). The variable \(Y_v\) is measurable without
coordinate \(i\); the conditional pullout identity therefore gives
\[
 E[Y_uY_v]
  =E\!\left[Y_v E[Y_u\mid\mathcal G_{I\setminus\{i\}}]\right]=0.
\]
Cauchy--Schwarz makes the product integrable. If one set is empty,
the same assertion also follows from the zero mean of the other component.
Squaring the finite sum and taking expectation proves the variance formula.

For uniqueness, let \(H_u\) be another decomposition with the stated
coordinate measurability and conditional zero-average conditions.
Averaging \(H_u\) onto \(\mathcal G_v\) gives \(H_u\) when \(u\subseteq v\)
and gives zero otherwise: in the latter case integrate a coordinate of
\(u\setminus v\) first and use its zero-average condition.
Consequently \(P_vY=\sum_{u\subseteq v}H_u\).
Finite subset inversion gives exactly the displayed formula for \(Y_v\),
so \(H_u=Y_u\) almost surely for every \(u\).
\end{proof}
\end{annexitem}

\begin{annexitem}{annex:p15.global.indices}
 {thm:p15-global-indices}{First, total and interaction contributions}
\begin{proof}
The decomposition gives
\(E[Y\mid X_i]=EY+Y_{\{i\}}\), hence its variance is \(V_{\{i\}}\).
Averaging over \(X_i\) removes exactly those components whose index sets
contain \(i\), so
\[
 E[Y\mid X_{-i}]=\sum_{u\subseteq I\setminus\{i\}}Y_u.
\]
Orthogonality and the law of total variance give
\[
 E[\operatorname{Var}(Y\mid X_{-i})]
    =V-\operatorname{Var}(E[Y\mid X_{-i}])
    =\sum_{u\ni i}V_u.
\]
All component variances are nonnegative. Dividing by \(V>0\) proves
the individual bounds and \(\sum_iS_i\leq1\). Also
\[
 \sum_iS_{T_i}=\frac{\sum_{\varnothing\ne u\subseteq I}|u|V_u}{V},
\]
which is a weighted average of integers between \(1\) and \(d\).
This proves the two total-sum bounds.

Equality \(\sum_iS_i=1\) holds exactly when every higher-order \(V_u\)
vanishes. Since \(V_u=E[Y_u^2]\), this is exactly the vanishing of all
interaction components almost surely, leaving an additive response.
Moreover
\(V(S_{T_i}-S_i)=\sum_{u\ni i,\ |u|\geq2}V_u\).
All these differences vanish precisely in the same additive case.

Finally the conditional-variance numerator is
\(E[(Y-E[Y\mid X_{-i}])^2]\).
It is zero exactly when \(Y=E[Y\mid X_{-i}]\) almost surely.
That is exactly measurability as a function of \(X_{-i}\), by the
measurable-factorization result for real coordinate tuples.
\end{proof}
\end{annexitem}

\begin{annexitem}{annex:p15.global.product-case}
 {thm:p15-global-product-case}{The centered bilinear counterexample}
\begin{proof}
Each uniform input has mean \(0\) and second moment
\(\frac12\int_{-1}^1x^2\,dx=1/3\).
Independence gives \(EY=0\) and \(E[Y^2]=1/9\).
Conditioning on either coordinate gives zero, so both first-order
numerators vanish. Conditioning on \(X_2\) gives conditional variance
\(X_2^2/3\); its expectation is \(1/9\).
The other total numerator is the same. Dividing by \(V=1/9\) yields
both total indices equal to \(1\).
The orthogonal decomposition therefore consists only of \(Y_{\{1,2\}}=Y\).
Every point on either coordinate axis has product zero.
\end{proof}
\end{annexitem}

\begin{annexitem}{annex:p15.global.pick-freeze}
 {thm:p15-global-pick-freeze}{Pairing identities and consistent estimates}
\begin{proof}
Write \(m_i(x_i)=E[f(X)\mid X_i=x_i]\), choosing the product-integral
version. Given \(A_i\), the remaining coordinates of \(A\) and \(B\)
are independent copies. Fubini or conditional independence gives
\[
 E[f(A)f(C_i)\mid A_i]=m_i(A_i)^2.
\]
Taking expectation and subtracting \((EY)^2\) proves the first identity.
The product is integrable by Cauchy--Schwarz, since each output has
the same square-integrable marginal law as \(Y\).

Given \(B_{-i}\), the two outputs \(f(B_i,B_{-i})\) and
\(f(A_i,B_{-i})\) are independent with the same conditional law.
Expanding their squared difference gives twice their conditional variance.
Taking expectation proves the second identity.

Across independent input pairs, both integrable paired quantities obey
the strong law. The sample mean of \(f(A)\) converges to \(EY\), and
that of \(f(A)^2\) converges to \(E[Y^2]\). Their difference-of-moments
variance estimate converges to \(V>0\), so it is positive eventually
almost surely. The paired numerators divided by this estimate converge
to the two indices. A ratio of random estimators is not thereby unbiased.
The needed integrable strong law is
Theorem~\ref{thm:p15-iid-series-strong-law}.

For each pair evaluate \(f(A)\), \(f(B)\), and \(f(C_i)\) for
\(i=1,\ldots,d\). These are \(d+2\) evaluations and can be shared
between the first and total computations. Multiplication by \(n\)
gives the stated nominal cost.
\end{proof}
\end{annexitem}

\begin{annexitem}{annex:p15.global.source-toys}
 {thm:p15-global-source-toys}{Product moments and the mixed bump contrast}
\begin{proof}
For \(X\) uniform on \([-1,1]\), let \(W=|4X-2|\).
The variable \(4X-2\) is uniform on \([-6,2]\), hence
\[
 EW=\frac18\left(\int_{-6}^0(-z)\,dz+\int_0^2z\,dz\right)=\frac52,
 \qquad E[W^2]=\frac{28}{3}.
\]
Thus \(\operatorname{Var}(W)=28/3-25/4=37/12\).
Translation and scaling give the stated \(m_i,v_i\).

Write \(g_i=(|4X_i-2|+a_i)/(1+a_i)\) and put \(h_i=g_i-m_i\),
so \(Eh_i=0\), \(Eh_i^2=v_i\), and the
\(h_i\) are independent. Expanding the product gives
\[
 G=\sum_{u\subseteq I}
       \left(\prod_{i\in u}h_i\right)
       \left(\prod_{j\notin u}m_j\right).
\]
Each nonempty component has zero average in each of its active
coordinates, so by uniqueness this is its functional decomposition.
Independence gives its squared norm
\(\prod_{i\in u}v_i\prod_{j\notin u}m_j^2\).
Summing these products over nonempty subsets gives
\(\prod_i(m_i^2+v_i)-\prod_i m_i^2\).
Summing only singleton \(i\), or all subsets containing \(i\),
gives the two numerator formulas.

Every term of \(D\) other than its rational bump is a function of
\(x_1\) alone or of \(x_2\) alone. Such terms cancel in the mixed
four-point contrast. The bump values at the four indicated points are
respectively \(3,\ 3/10,\ 3/5,\ 3/14\).
Therefore the remaining contrast is
\(3-3/10-3/5+3/14=81/35\).
\end{proof}
\end{annexitem}

\begin{annexitem}{annex:p15.global.dependent-boundary}
 {thm:p15-global-dependent-boundary}{The perfectly dependent example}
\begin{proof}
For either conditioning sigma-algebra, total variance gives a
nonnegative explained variance no larger than \(V\), and a nonnegative
mean conditional variance no larger than \(V\). Thus each ratio is
in \([0,1]\), without an independence assumption.

In the example, \(V=\operatorname{Var}(U)=1/4\).
Either \(X_1\) or \(X_2\) determines \(Y=U\), so both conditional
means equal \(Y\) and both conditional variances given the other input
are zero. The four ratios are consequently \(1,1,0,0\).
Independent resampling of the two fair marginals instead gives positive
probability to \((0,1)\) and \((1,0)\), points absent from the original
diagonal joint support. It therefore changes the model-input law.
\end{proof}
\end{annexitem}

\begin{annexitem}{annex:p15.global.causal-boundary}
 {thm:p15-global-causal-boundary}{Equal observational sensitivities, different interventions}
\begin{proof}
Let \(U_1,U_2\) be independent fair binary variables. In both models set
\(X_1=U_1\) and \(X_2=U_2\).
In model A set \(Y=X_1\); in model B set \(Y=U_1\).
Observationally both yield \(Y=X_1\), with the same complete joint law
of the three observed variables. In particular the observed inputs are
independent and the same model-output indices are obtained in both:
factor \(1\) has first and total index \(1\), and factor \(2\) has
both indices zero.

An intervention setting \(X_1=1\) replaces only that structural equation.
In model A it makes \(Y=1\) surely. In model B, \(Y=U_1\) remains a
fair binary variable. Their post-intervention expectations are \(1\)
and \(1/2\), respectively.
The observational law and its sensitivity indices therefore do not
identify the intervention response, even with independent observed inputs.
\end{proof}
\end{annexitem}

\begin{annexitem}{annex:p15.global.ranking}
 {thm:p15-global-ranking}{Input ranges control the linear ranking}
\begin{proof}
Independence and centering give
\(V=v_1+4v_2\).
The model is additive, with component variances \(v_1\) and \(4v_2\),
so the index formulas follow.
A uniform variable on \([-a,a]\) has variance \(a^2/3\), by direct
integration. For both half-widths \(1\), the two contributions are
\(1/3\) and \(4/3\), yielding \(1/5,4/5\).
With the first half-width \(3\), they become \(3\) and \(4/3\),
yielding \(9/13,4/13\).

The maximum of \(x_1+2x_2\) on the first rectangle is \(3\);
on the second it is \(5\). Thus the asserted threshold changes status.
In fact the second law has
\[
 P(Y>3)=\frac1{12}\int_0^1
                   \bigl(3-(3-2x_2)\bigr)\,dx_2=\frac1{12}.
\]
The event is therefore not confined to a measure-zero corner.

For the variance family take
\(0\leq\ell_i\leq u_i\) and exclude cases with zero total variance.
Since the two index denominators are the same positive number,
\(S_1\geq S_2\) is equivalent to \(v_1\geq4v_2\).
It holds on the whole independently varying rectangle exactly when
its smallest left side \(\ell_1\) is at least its largest right side
\(4u_2\). If the inequality fails, the endpoint pair
\((\ell_1,u_2)\) has positive output variance and violates the ranking.
\end{proof}
\end{annexitem}
